% Zeta-преобразователь.
% Два дросселя L1, L2, связующий конденсатор C1, ключ VT (MOSFET с
% индуцированным каналом), диод VD1, выходной C и R_н.
% U_вых = U_вх·D/(1-D); в отличие от Sepic, L2 стоит на выходе —
% выходной ток непрерывен.
\documentclass[border=4pt]{standalone}
\usepackage{amsmath}
\usepackage{fontspec}
\setmainfont{Calibri}
\usepackage[european, RPvoltages]{circuitikz}
\usetikzlibrary{arrows.meta, calc}
\definecolor{ink}{RGB}{31,42,51}
\definecolor{grp}{RGB}{46,110,140}
\ctikzset{bipoles/length=1.0cm}
\ctikzset{bipoles/thickness=0.9}
\ctikzset{diodes/fill=ink}
\ctikzset{inductor=american}

\begin{document}
\begin{tikzpicture}[line width=1.0pt, color=ink, font=\small]
  % --- входные зажимы (открытые)
  \draw (0.6,3.6) node[ocirc]{} node[above left]{$+$} -- (1.1,3.6);
  \node[anchor=east] at (0.4,1.8) {$U_{\text{вх}}$};
  % --- L1 в шине входа
  \draw (1.1,3.6) to[L, l_=$L_1$] (2.5,3.6) -- (3.1,3.6) node[circ]{};
  % --- диод VD1: анод на земле, катод в узле A (остриё вверх)
  \draw (3.1,0) node[circ]{} to[D, name=VD1] (3.1,3.6);
  \node[anchor=west] at (3.4,1.8) {VD1};
  % --- связующий конденсатор C1 (горизонтально)
  \draw (3.1,3.6) -- (3.5,3.6) to[C, l_=$C_1$] (4.7,3.6) -- (5.1,3.6)
        node[circ]{};
  % --- ключ VT: земля -> узел B (путь снизу вверх, затвор влево)
  \draw (5.1,0) node[circ]{} to[Tnigfete, name=VT] (5.1,2.8)
        -- (5.1,3.6);
  \node[anchor=west] at (5.6,2.0) {VT};
  \draw[grp] (VT.G) -- ++(-0.5,0);
  \node[anchor=east, text=grp] at ($(VT.G)+(-0.55,0)$) {ШИМ};
  % --- L2: узел B -> выход (выходной ток непрерывен)
  \draw (5.1,3.6) -- (5.5,3.6) to[L, l_=$L_2$] (6.9,3.6) -- (8.9,3.6);
  % --- нижняя шина
  \draw (0.6,0) node[ocirc]{} node[below left]{$-$} -- (8.9,0);
  % --- выходной конденсатор и нагрузка
  \draw (7.5,3.6) node[circ]{} to[C, l_=$C$] (7.5,0) node[circ]{};
  \draw (8.3,3.6) node[circ]{} to[R, l=$R_{\text{н}}$] (8.3,0)
        node[circ]{};
  % --- выходные зажимы
  \draw (8.9,3.6) node[ocirc]{} node[above right]{$+$};
  \draw (8.9,0) node[ocirc]{} node[below right]{$-$};
  \node[anchor=west] at (9.1,1.8) {$U_{\text{вых}}$};
\end{tikzpicture}
\end{document}
